\scuAbstractEnHeading
Polymer conformation connects molecular structure with material response. This illustrative thesis starts from the freely jointed chain and derives the mean-square end-to-end distance by separating diagonal and cross terms in the segment-vector sum. The vanishing cross terms follow from the assumed independence and isotropy of the segments. Effective segment length is distinguished from the geometric length of a chemical repeat unit.

Ideal-chain behavior is then compared with a scaling description of self-avoiding chains in a good solvent. The interpretation of a fitted exponent depends on chain length, solvent quality, spatial dimension, and the observation scale. A straight line on logarithmic axes alone does not establish an asymptotic scaling regime. Dimensionless curves calculated from stated formulas illustrate these distinctions; they are neither measurements nor predictions of a particular material.

The analysis emphasizes explicit model assumptions, clearly defined observables, and sensitivity checks across chain-length ranges. Experimental validation would additionally require independently characterized molecular-weight distributions and solvent conditions. This document is a theoretical and typesetting example without original experimental data. Its content must be replaced by the author's own research before submission.

\par\vspace{20bp}\noindent\textbf{Keywords: }\scuKeywordsEn
